% !TEX TS-program = XeLaTeX
% !TEX encoding = UTF-8 Unicode

\chapter*{\hfill 引　　言 \hfill}
\defaultfont
\addcontentsline{toc}{chapter}{引　　言}
\label{chap:intro}

脊柱健康与人体长期姿态、运动习惯和日常生活方式密切相关。脊柱后凸、脊柱侧弯等姿态异常在早期往往表现为身体形态和运动姿态的细微差异，若能够通过低成本、非接触的方式进行持续感知，可为日常健康监测和早期辅助筛查提供技术支持。传统姿态检测方法多依赖摄像头、可穿戴设备或专业医学检测设备。摄像头方案容易受到视角、光照和隐私问题影响；可穿戴设备需要用户主动佩戴，长期使用的舒适性和依从性有限；专业医学设备检测精度高，但部署成本较高，难以覆盖普通室内场景下的连续感知需求。

近年来，基于 WiFi 信号的人体感知技术受到广泛关注。WiFi 设备在室内环境中普遍存在，无线信号在传播过程中会受到人体反射、遮挡、散射和多径效应影响。信道状态信息（Channel State Information，CSI）能够描述不同子载波和天线链路上的细粒度信道响应。相较于接收信号强度（Received Signal Strength Indicator，RSSI），CSI 同时包含幅度和相位信息，能够更细致地反映人体姿态对无线信道的扰动。因此，利用 CSI 进行非接触式人体姿态检测具有较好的研究价值和应用潜力。

与行走、跌倒和手势等具有明显动态过程的行为识别任务相比，脊柱姿态感知更加关注人体躯干形态及其相对稳定状态所造成的信道差异。这类差异通常叠加在环境多径、硬件相位偏移、人员个体差异和随机噪声之上，难以通过单一时刻或单一链路的观测直接判断。因此，研究的关键不只是选择分类模型，还包括如何构造对姿态变化敏感、对公共干扰具有一定抑制能力的信号表示，以及如何建立与实际检测单位一致的训练和评价流程。

从研究目标看，基于无线信号的姿态检测属于间接感知问题。系统观测到的是人体对传播信道造成的扰动，而不是人体骨骼或脊柱的解剖结构。因而，本文所讨论的检测结果用于描述无线感知条件下的姿态类别，不等同于临床诊断结论。明确这一边界有助于合理解释实验指标，也有助于避免将特定数据条件下的分类性能外推到未经验证的人群和环境。

本文面向脊柱姿态检测任务，研究一种基于 WiFi CSI 的三分类识别方法。该方法将检测目标划分为正常姿态、脊柱后凸相关姿态和脊柱侧弯相关姿态。系统首先解析原始 CSI 文件，获得时间帧、子载波、接收天线和发射天线维度上的复数信道响应；随后构造接收天线相位比特征，并沿时间维度进行相位展开；在此基础上通过滑动窗口组织 CSI 序列，再将每个原始文件的窗口特征聚合为文件级统计特征；最后使用 ExtraTrees 集成分类器完成姿态类别预测。

围绕上述任务，本文重点研究以下三个问题：第一，接收天线之间的相对相位关系能否削弱公共相位误差的影响，并保留与人体姿态有关的空间链路差异；第二，文件级统计聚合能否将连续且长度可变的 CSI 序列转化为相对稳定的固定长度表示；第三，在样本规模有限、特征维度较高的条件下，集成学习方法能否获得稳定的文件级分类性能。三个问题分别对应信号表示、样本组织和分类决策三个层次，共同构成本文的方法框架。

本文的主要工作包括以下几个方面：
\begin{enumerate}
\item 分析 WiFi CSI 用于脊柱姿态检测的可行性，说明人体姿态变化与无线信道幅度、相位扰动之间的关系。
\item 构建基于接收天线相位比的 CSI 特征处理方法，削弱硬件相位偏移和公共相位误差对分类特征的影响。
\item 设计文件级统计特征聚合流程，将连续 CSI 相位比序列转换为固定长度特征向量，便于传统机器学习模型进行分类。
\item 采用 ExtraTrees 集成分类器建立三类脊柱姿态检测模型，并与 TCN、SVM 和 Random Forest 等方案进行比较。
\item 在文件级留出测试集上进行实验验证，分析准确率、精确率、召回率、F1 值和混淆矩阵结果。
\end{enumerate}

全文结构安排如下：第一章介绍相关研究与技术基础；第二章说明 CSI 数据表示与特征处理方法；第三章给出基于相位比统计特征和 ExtraTrees 的脊柱姿态检测方法；第四章介绍实验设计、模型对比和测试结果；第五章阐述系统实现与检测流程；第六章讨论方法适用条件、泛化问题、潜在混杂因素和应用边界；最后对全文工作进行总结。
